%\textbf{\emph{ejemplo}}
\input{Datos.txt}
\documentclass[letterpaper,9pt]{article}
\usepackage{ifthen}

% *** Idioma ***
\usepackage[spanish,es-nodecimaldot]{babel}
\usepackage[utf8]{inputenc}

% *** Fuente ***
\usepackage{fontspec}
\setmainfont{Noto Sans}


% *** Colores ***
\usepackage{xcolor}
\definecolor{dorado}{RGB}{190,150,91} 
\definecolor{gris}{rgb}{0.752,0.752,0.752}

% *** Margenes y Tamaño de Hoja
\usepackage{geometry}
\geometry{
left = 15mm,
right = 15mm,
top = 45mm,
bottom = 10mm,
headheight = 20mm
}

% *** Imagenes ***
\usepackage{graphicx}
\graphicspath{./imagenes}
\usepackage{float}
\usepackage{subcaption}
\usepackage{wrapfig}
\usepackage{tabularray}
%OJO con esta línea
\UseTblrLibrary{booktabs}

% Quita las leyenda de longtblr
\DefTblrTemplate{contfoot-text}{default}{}
\DefTblrTemplate{conthead-text}{default}{}
\DefTblrTemplate{capcont}{default}{}

% *** Fecha ***
%\usepackage[useregional]{datetime2}
\usepackage{datetime2}
\DTMsetdatestyle{ddmmyyyy}

%No cortar palabras
\usepackage[none]{hyphenat}


% *** Encabezado ***
\usepackage{lastpage} 
\usepackage{fancyhdr}
\pagestyle{fancy}
\renewcommand{\headrulewidth}{0pt}

\lhead{
		\begin{tblr}{
		width = \linewidth,
		colspec = {Q[120]Q[507]Q[70]Q[70]},
		row{2} = {c},
		row{3} = {c},
		column{2} = {c},
		cell{1}{1} = {r=4}{c},
		cell{1}{2} = {r=3}{},
		cell{4}{2} = {h},
		cell{1}{3} = {c},
		cell{1}{4} = {c},
		cell{4}{3} = {c},
		cell{4}{4} = {c},
		vlines,
                      hlines,
	}
	\includegraphics[width=0.15\textwidth]{imagenes/Logo_de_la_Conagua} & {\textbf{COMISIÓN NACIONAL DEL AGUA.}\\Gerencia de Tecnologías de Información y Comunicaciones.\\ Subgerencia de Soporte Técnico,  Telecomunicaciones y Seguridad.} & \footnotesize Página   & {\footnotesize \thepage\ de \pageref{LastPage}} \\
	&                                                                                       & \footnotesize Versión   & \footnotesize Pendiente \\
	&                                                                                       & \footnotesize Fecha & {\footnotesize Pendiente }    \\
	&\ { \textbf{ Solicitud de Ampliación del Servicio de Internet }}                                       &  \footnotesize No. de formato      &  \footnotesize \NOFORMATO        
\end{tblr}
}
\chead{}
\rhead{}

% *** Pie de Pagina
\lfoot{}
\cfoot{}

\begin{document}
\sloppy

\begin{longtblr}[
	label = none,
	entry = none,
	]{
		width = \linewidth,
		colspec = {Q[1]Q[1.6]Q[1]Q[1.6]Q[2]Q[1]},  
                    %  row{1}={r}   
 		cell{1}{2}={c},
 		cell{1}{4}={c},
 		cell{1}{6}={c},
	}
{\textbf{Memorando:}}&\NOMEMO&{\textbf{No. Ticket:}}&\NOTICKET&{\textbf{Fecha de la solicitud:}} &\FECHASOLI    
\end{longtblr}

\vspace{-30pt}
\begin{longtblr}[
	label = none,
	entry = none,
	]{
		width = \linewidth,
		colspec = {Q[200]Q[300]Q[90]Q[90]},
		row{1} = {c},               
		cell{1}{1} = {c=4}{0.874\linewidth},
		cell{2}{2} = {c=3}{0.674\linewidth},
		cell{3}{2} = {c=3}{0.674\linewidth},
		cell{4}{2} = {c=3}{0.674\linewidth},
		cell{5}{2} = {c=3}{0.674\linewidth},     
		cell{6}{2} = {c=3}{0.674\linewidth},     
		cell{7}{2} = {c=3}{0.674\linewidth},     
		%cell{8}{2} = {c=3}{0.674\linewidth},     
		hlines,
		vlines,
	}
\textbf{Datos de Usuario (a)} &  &      &  \\
Nombre:                     &  \NOMBREUSUARIO \\
Puesto o Cargo: & \PUESTOUSUARIO\\
Unidad Aministrativa:           &\UAUSUARIO\\
Área:&\AREAUSUARIO\\
Correo: &   \CORREOUSUARIO  \\
Dirección:   &  \textbf  \DIRECCION  \\
Teléfono: &\TELUSUARIO &Extensión:&\EXTUSUARIO
\end{longtblr}
\newboolean{ipactual}
\setboolean{ipactual}{\IPACTUAL}
\vspace{-37pt}
\ifthenelse{\boolean{ipactual}}{
\begin{longtblr}[
	label = none,
	entry = none,
	]{
		width = \linewidth,
		colspec = {Q[200]Q[300]Q[90]Q[90]},          
		cell{1}{2} = {c=3}{0.674\linewidth}, 
		hlines,
		vlines,
	}
IP del equipo asignado: & \IPUSUARIO&&
\end{longtblr}
}{
\begin{longtblr}[
	label = none,
	entry = none,
	]{
		width = \linewidth,
		colspec = {Q[200]Q[300]Q[90]Q[90]},	 
		cell{1}{2} = {c=3}{0.674\linewidth},    
		cell{2}{2} = {c=3}{0.674\linewidth},      
		hlines,
		vlines,
	}
IP anterior del equipo: & \IPANTERIOR&&\\
IP actual del equipo asignado: & \IPUSUARIO
\end{longtblr}
}

\vspace{-30pt}
\begin{longtblr}[
	label = none,
	entry = none,
	]{
		width = \linewidth,
		colspec = {Q[200]Q[200]Q[200]},                   
		cell{1}{1} = {c=3}{0.974\linewidth},
		cell{3}{1} = {c=3}{0.974\linewidth},
	}
{\textbf {Perfil básico:}} Contempla las siguientes categorías, necesarias para el desempeño sustantivo y la adecuada gestión de los recursos informáticos &  &     \\
{•	Negocios y economía\\
•	Información sobre computadoras \\ \hspace*{2mm} e internet\\
•	Instituciones educativas\\
•	Entretenimiento y artes\\
•	Gobierno\\}
               & 
{•	Salud y medicina\\
•	Portales de internet\\
•	Legal\\
•	Noticias\\
•	Bienes raíces\\}
                &
{• Motores de búsqueda\\
•	Educación sexual\\
•	Compras\\
•	Viajar\\
• Correo electrónico basado \\ \hspace*{3mm}en la web.} \\

{\textbf {Perfil comunicación social:}} Incorpora el perfil básico y contempla opciones adicionales que deberá selecionar y justificar plenamente, autorizadas por el Director o Gerente del área de adscripción.
\end{longtblr}

\vspace{-30pt}
\begin{longtblr}[
	label = none,
	entry = none,
	]{
		width = \linewidth,
		colspec = {Q[18]Q[200]Q[18]Q[200]Q[18]Q[200]},
		row{1} = {c},
		cell {2}{6}={l},
		cell{1}{1} = {c=6}{0.874\linewidth},	
		hlines,
		vlines,
	}
\textbf{Opciones de Navegación}         \\
 (\ALMACENAMIENTO )& \textbf{Almacenamiento y copia de seguridad en línea}  & (\BLOGS )& \textbf{Sitios personales y blogs} & (\SHAREWARE)& \textbf{Descarga de software/drivers}  \\
(\REDES)& \textbf{Redes sociales}  & (\TRANSMISION )& \textbf{Transmisión de medios} & (\OTRA)& \textbf{Otra}  
\end{longtblr}
\vspace{-5mm}

{\textbf{\\Importante:}} En la(s) siguiente(s) tabla(s) deberá JUSTIFICAR cada una de las opciones seleccionadas; considerando en lo particular las solicitudes deben estar asociadas a requerimientos específicos de las funciones encomendadas e información institucional y NO para fines ni servicios personales. 

%%%JUSTIFICACIONES
\newboolean{almacenamiento}
\setboolean{almacenamiento}{\ALMACENAMIENTOBOOL}
\newboolean{blogs}
\setboolean{blogs}{\BLOGSBOOL}
\newboolean{shareware}
\setboolean{shareware}{\SHAREWAREBOOL}
\newboolean{redes}
\setboolean{redes}{\REDESBOOL}
\newboolean{transmision}
\setboolean{transmision}{\TRANSMISIONBOOL}
\newboolean{otra}
\setboolean{otra}{\OTRABOOL}
\newboolean{otra2}
\setboolean{otra2}{\OTRABOOLdos}
\newboolean{otra3}
\setboolean{otra3}{\OTRABOOLtres}
\newboolean{otra4}
\setboolean{otra4}{\OTRABOOLcuatro}

\vspace{-3mm}
\textbf {\\JUSTIFICACIONES\\}
%%%ALMACENAMIENTO
\ifthenelse{\boolean{almacenamiento}}{
\vspace{2mm}
\noindent
\begin{tblr}[
	label = none,
	entry = none,
	]{
		width = \linewidth,
		colspec = {Q[150]Q[300]},              
		cell{1}{1} = {c=2}{0.874\linewidth},  
		hlines,
		vlines,
	}
\textbf{Almacenamiento y copia de seguridad en línea} \\
\textbf{Justificación:}&
\JUSTIFICAALMACENAMIENTO
\end{tblr} 
}{~}
%%BLOGS
\ifthenelse{\boolean{blogs}}{
\vspace{2mm}
\noindent
\begin{tblr}[
	label = none,
	entry = none,
	]{
		width = \linewidth,
		colspec = {Q[150]Q[300]},             
		cell{1}{1} = {c=2}{0.874\linewidth},
		hlines,
		vlines,
	}
\textbf{Sitios personales y blogs} \\
\textbf{Justificación:}&
\JUSTIFICABLOGS
\end{tblr}
}{~}
%%SHAREWARE
\ifthenelse{\boolean{shareware}}{
\vspace{2mm}
\noindent
\begin{tblr}[
	label = none,
	entry = none,
	]{
		width = \linewidth,
		colspec = {Q[150]Q[300]},                  
		cell{1}{1} = {c=2}{0.874\linewidth}, 
		hlines,
		vlines,
	}
\textbf{Descarga de software / drivers} \\
\textbf{Justificación:}&
\JUSTIFICASHAREWARE
\end{tblr}
}{~}
%%%REDES
\ifthenelse{\boolean{redes}}{
\vspace{2mm}
\noindent
\begin{tblr}[
	label = none,
	entry = none,
	]{
		width = \linewidth,
		colspec = {Q[150]Q[300]},              
		cell{1}{1} = {c=2}{0.874\linewidth}, 
		hlines,
		vlines,
	}
\textbf{Redes Sociales} \\
\textbf{Justificación:}&
\JUSTIFICAREDES
\end{tblr}
}{~}

%%TRANSMISION DE MEDIOS
\ifthenelse{\boolean{transmision}}{
\vspace{2mm}
\noindent
\begin{tblr}[
	label = none,
	entry = none,
	]{
		width = \linewidth,
		colspec = {Q[150]Q[300]},                 
		cell{1}{1} = {c=2}{0.874\linewidth},
		hlines,
		vlines,
	}
\textbf{Transmisión de medios} \\
\textbf{Justificación:}&
\JUSTIFICATRANSMISION
\end{tblr}
}{~}

%%OTRA
\ifthenelse{\boolean{otra}}{
\vspace*{3mm}
\noindent
\begin{tblr}[
	label = none,
	entry = none,
	]{
		width = \linewidth,
		colspec = {Q[150]Q[300]},               
		cell{1}{1} = {c=2}{0.874\linewidth},   
		hlines,
		vlines,
	}
\textbf{Otra: \OTRAC} \\
{Referencia del servicio requerido\\ (URL del servicio):} & \URLOTRA\\
\textbf{Justificación:}&
\JUSTIFICAOTRA
\end{tblr}
}{~}
%%otra2
\ifthenelse{\boolean{otra2}}{
\par
\vspace*{3mm}
\noindent
\begin{tblr}[
	label = none,
	entry = none,
	]{
		width = \linewidth,
		colspec = {Q[150]Q[300]},                   
		cell{1}{1} = {c=2}{0.874\linewidth},   
		hlines,
		vlines,
	}
\textbf{Otra: \OTRACdos} \\
{Referencia del servicio requerido\\ (URL del servicio):} & \URLOTRAdos\\
\textbf{Justificación:}&
\JUSTIFICAOTRAdos
\end{tblr}
}{~}
%%otra3
\ifthenelse{\boolean{otra3}}{
\par
\vspace*{3mm}
\noindent
\begin{tblr}[
	label = none,
	entry = none,
	]{
		width = \linewidth,
		colspec = {Q[150]Q[300]},              
		cell{1}{1} = {c=2}{0.874\linewidth},   
		hlines,
		vlines,
	}
\textbf{Otra: \OTRACtres} \\
{Referencia del servicio requerido\\ (URL del servicio):} & \URLOTRAtres\\
\textbf{Justificación:}&
\JUSTIFICAOTRAtres
\end{tblr}
}{~}
%%otra4
\ifthenelse{\boolean{otra4}}{
\par
\vspace*{3mm}
\noindent
\begin{tblr}[
	label = none,
	entry = none,
	]{
		width = \linewidth,
		colspec = {Q[150]Q[300]},      
		cell{1}{1} = {c=2}{0.874\linewidth},   
		hlines,
		vlines,
	}
\textbf{Otra: \OTRACcuatro} \\
{Referencia del servicio requerido\\ (URL del servicio):} & \URLOTRAcuatro\\
\textbf{Justificación:}&
\JUSTIFICAOTRAcuatro
\end{tblr}
}{~}

\clearpage
%%%..... POLITICAS
\section*{Términos y Condiciones del servicio:}
{\small \begin{enumerate}
	\item 	El formato deberá estar debidamente llenado y contener toda la información requerida facilitando la aplicación expedita de las configuraciones solicitadas.
	\item 	El solicitante deberá presentar este formato adjuntando el memorando y número de caso (ticket) de Mesa de Servicio asociado, sin los cuales no se podrá atender su solicitud.         
	\item  	El solicitante deberá proporcionar la dirección IP de su equipo de cómputo para poder aplicar las configuraciones requeridas. Esta información la podrá obtener utilizando el comando “ipconfig /all” en una ventana de línea de comando. En caso de requerir ayuda para ejecutar el comando indicado, favor de contactar al área de Soporte Técnico de la CONAGUA.
	\item 	El solicitante deberá agregar en la \textbf {Justificación} si la solicitud se deriva de un cambio de lugar (oficina, mampara o piso) del usuario que a su vez haya derivado en un cambio de dirección IP de su equipo de cómputo.
	\item  	En caso de un cambio de dirección IP, el usuario deberá especificar la dirección IP anterior para eliminar los privilegios en dicha dirección IP. Si el solicitante NO indica que se trata de un cambio de dirección IP, éste será responsable de cualquier acceso no autorizado que se derive de los permisos de la dirección IP anterior al no tramitar la baja correspondiente.
	\item 	El usuario es responsable del uso que se otorga con el acceso ampliado otorgado a su equipo de cómputo, por lo que deberá vigilar que el uso sea acorde a las Políticas de la Seguridad de la Información definidas por la Gerencia de Tecnología de la Información y Comunicaciones.
	\item 	El solicitante deberá conservar el Acuse o copia del formato firmado y sellado, así como el memorando asociado, para posteriores aclaraciones.
	\item  	Es responsabilidad de los gerentes y subgerentes llevar un control de los usuarios y sus direcciones IP’s con accesos de Internet ampliados.
	\item 	Al firmar el usuario se da por enterado de las Políticas del servicio y acepta la responsabilidad de cualquier uso inadecuado que se le dé a los privilegios de acceso ampliados los cuales haya solicitado.
	\item 	Al firmar el Gerente o Director que autoriza se da por enterado de las Políticas del servicio y acepta la corresponsabilidad del uso que le dé el usuario al acceso ampliado otorgado.
\end{enumerate}
%%nombres y firmas
\begin{longtblr}[
	label = none,
	entry = none,
	]{
			width = \linewidth,
			colspec = {Q[100]Q[100]Q[100]},
			row{1} = {c},
                                row{2}={c},
		}
		%Fecha: \today                      &                                                                                                       \\
		{Usuario que acepta\\~\\[1cm] \rule{5cm}{0.5mm}\\ \NOMBREUSUARIO \\~\\ \PUESTOUSUARIO} &
                     {Gerente o Director Local que Autoriza\\~\\[1cm] \rule{5cm}{0.5mm}\\ \NOMBREJEFE \\~ \\ \PUESTOJEFE}
	\end{longtblr}
{\textbf{Nota:}}
Este formato debe ser entregado a través del servicio de Mesa de Servicio, acompañado de un Memorando dirigido al {\textbf {Gerente de Tecnologías de la Información y Comunicaciones}} con copia a {\textbf {Subgerente de Soporte Técnico, Telecomunicaciones y Seguridad}}.
De no recibir dichos formatos en original en un lapso no mayor a 5 días hábiles se retirarán los permisos solicitados y se otorgará el acceso a navegación básico.

\end{document}